\chapter{\texorpdfstring{致\quad 谢}{致谢}} %这里多出来的代码是用来消除hyperref包引入的警告的，\texorpdfstring{}{}命令的作用是告诉hyperref包在生成PDF书签时使用第二个参数，而不是第一个参数中的特殊字符。

时光荏苒，四年的本科生活即将画上句号。回首在天津大学管理与经济学部度过的点滴时光，心中充满了感激与不舍。在此论文付梓之际，谨向所有给予我关心、帮助和支持的师长、同学、朋友和家人表达最诚挚的谢意。

首先，我要向我的指导教师肖应钊老师致以最深切的谢意。从论文选题、文献研读、研究框架的搭建，到问卷设计、数据分析与论文写作的每一个环节，肖老师都给予了悉心的指导与耐心的修改建议。肖老师严谨的治学态度、深厚的学术造诣以及精益求精的工作作风，是我学习的榜样，将激励我在今后的人生道路上继续努力前行。

其次，感谢天津大学管理与经济学部的各位老师，是他们在专业课程上的悉心讲授让我夯实了管理学、心理学、统计学等多学科基础知识，为本论文的研究工作奠定了重要的理论基础。同时，感谢创业学院的老师们和同学们在问卷调查阶段给予的大力支持与配合，使本研究的实证数据得以顺利收集。

再者，感谢一同走过四年时光的同窗好友，与你们一起讨论、互相鼓励、共同成长的日子是我大学生涯里最珍贵的回忆。

最后，最深的感激献给我的家人。感谢父母多年来的悉心呵护与无私付出，你们的支持与理解是我不断前行的最大动力。

由于本人学识水平所限，论文中难免存在疏漏与不足之处，敬请各位专家、老师不吝指正。

\clearpage
